% =========================================================
% Lifting and Co-Lifting
% =========================================================

\section{Lifting and Co-Lifting}
\label{sec:lifting}

\begin{remark*}[Toolkit: Lifting and Co-Lifting]
A poset that lacks a bottom or top element can be given one by adjoining
a new universal bound.  Lifting adjoins a bottom; co-lifting adjoins a
top.  These constructions are dual to each other and are used throughout
domain theory, where a bottom element represents an undefined state.
\end{remark*}

% ---------------------------------------------------------
% def:lifting
% ---------------------------------------------------------

\begin{definitionbox}{Definition (Lifting)}
\begin{definition}[Lifting]
\label{def:lifting}
Let $(P, \leq)$ be a partially ordered set.  The \emph{lifting} of $P$
is the poset $P_{\bot} \coloneqq P \cup \{\bot\}$, where $\bot \notin P$,
with the order:
\[
\bot \leq x \quad \text{for all } x \in P_{\bot},
\qquad
x \leq_{P_{\bot}} y \;\Longleftrightarrow\; x \leq_P y
\quad \text{for all } x, y \in P.
\]
Thus $P_{\bot}$ is $P$ with a new
\hyperref[def:bottom-element]{bottom element} $\bot$ adjoined.
\end{definition}
\end{definitionbox}

\begin{remark*}[Standard quantified statement]
\[
P_{\bot} = P \cup \{\bot\},\quad
\bot \notin P,\quad
\forall x \in P_{\bot},\; \bot \leq x,\quad
\forall x,y \in P,\; x \leq_{P_\bot} y \Leftrightarrow x \leq_P y.
\]
\end{remark*}

\begin{remark*}[Interpretation]
Lifting adjoins a universal lower bound to a poset that may not have one.
In domain theory $\bot$ represents an undefined, divergent, or
unspecified state.  Any set $S$ viewed as an antichain can be lifted to
the \emph{flat poset} $S_{\bot}$, in which all elements of $S$ are
incomparable above $\bot$.
\end{remark*}

\begin{dependencies}
\hyperref[def:lattice-order-partial-order]{Partial order},
\hyperref[def:union]{Union},
\hyperref[def:bottom-element]{bottom element}.
\end{dependencies}

% ---------------------------------------------------------
% def:co-lifting
% ---------------------------------------------------------

\begin{definitionbox}{Definition}
\begin{definition}[Co-Lifting]
\label{def:co-lifting}
Let $(P, \leq)$ be a partially ordered set.  The \emph{co-lifting} of $P$
is the poset $P^{\top} \coloneqq P \cup \{\top\}$, where $\top \notin P$,
with the order:
\[
x \leq \top \quad \text{for all } x \in P^{\top},
\qquad
x \leq_{P^{\top}} y \;\Longleftrightarrow\; x \leq_P y
\quad \text{for all } x, y \in P.
\]
Thus $P^{\top}$ is $P$ with a new
\hyperref[def:top-element]{top element} $\top$ adjoined.
\end{definition}
\end{definitionbox}

\begin{remark*}[Standard quantified statement]
\[
P^{\top} = P \cup \{\top\},\quad
\top \notin P,\quad
\forall x \in P^{\top},\; x \leq \top,\quad
\forall x,y \in P,\; x \leq_{P^\top} y \Leftrightarrow x \leq_P y.
\]
\end{remark*}

\begin{remark*}[Interpretation]
Dual to lifting via the
\hyperref[prop:duality-principle]{Duality Principle}: co-lifting adjoins
a universal upper bound.  The co-lifting of an antichain produces a flat
poset in which all elements are incomparable below $\top$.
\end{remark*}

\begin{dependencies}
\hyperref[def:lifting]{Lifting},
\hyperref[def:top-element]{top element},
\hyperref[prop:duality-principle]{duality principle}.
\end{dependencies}
